\documentclass[12pt,a4paper]{report}

% ------------------ PACKAGES ------------------
\usepackage[a4paper,left=20mm,right=10mm,top=20mm,bottom=20mm]{geometry}
\usepackage{setspace}
\usepackage{titlesec}
\usepackage{indentfirst}
\usepackage{fancyhdr}
\usepackage{graphicx}
\usepackage{caption}
\usepackage{amsmath}
\usepackage{chngcntr}
\usepackage{enumitem}
\usepackage{float}
\usepackage{booktabs}
\usepackage{longtable}
\usepackage{listings}
\usepackage{newtxtext,newtxmath}

% ------------------ SPACING ------------------
\onehalfspacing
\setlength{\parindent}{1.25cm}
\setlength{\parskip}{0pt}

% ------------------ PAGE NUMBERING ------------------
\pagestyle{fancy}
\fancyhf{}
\fancyfoot[C]{\thepage}
\renewcommand{\headrulewidth}{0pt}

% ------------------ CHAPTER FORMAT ------------------
\titleformat{\chapter}
{\bfseries\centering\fontsize{13pt}{13pt}\selectfont}
{\thechapter}
{1em}
{}
\titlespacing*{\chapter}{0pt}{0pt}{1\baselineskip}

\titleformat{\section}
{\bfseries\fontsize{12pt}{12pt}\selectfont}
{\thesection}
{1em}
{}

% ------------------ NUMBERING ------------------
\numberwithin{equation}{chapter}
\numberwithin{figure}{chapter}
\numberwithin{table}{chapter}

\captionsetup[figure]{labelfont=bf, justification=centering}
\captionsetup[table]{font=small}

% ------------------ ENUMERATION STYLE ------------------
\setlist[itemize]{label=-, leftmargin=1.25cm}
\setlist[enumerate,1]{label=\alph*)}
\setlist[enumerate,2]{label=\arabic*)}

% ------------------ CODE STYLE ------------------
\lstset{
    basicstyle=\ttfamily\fontsize{9pt}{9pt}\selectfont,
    breaklines=true,
    numbers=left,
    numberstyle=\tiny,
    frame=single,
    tabsize=2
}

% ------------------ DOCUMENT ------------------
\begin{document}

% ------------------ TITLE PAGE ------------------
\begin{titlepage}
\centering
\begin{figure}[htbp]
\centering
\IfFileExists{img/utm_logo.png}{\includegraphics[width=0.2\textwidth]{img/utm_logo.png}}{\vspace{2cm}}
\end{figure}
MINISTRY OF EDUCATION AND RESEARCH\\
OF THE REPUBLIC OF MOLDOVA\\
Technical University of Moldova\\
Faculty of Computers, Informatics and Microelectronics\\
Department of Software Engineering and Automation\\[2cm]

Vasilean Vitalie\\
Group FAF-241\\[1.5cm]

\textbf{\Large Report}\\
Laboratory Work No. 9\\
Operating Systems\\
Topic: Watching System Calls with strace (Track A - Linux)\\[2cm]

\begin{flushright}
Checked by:\\
univ. assist. / lect. Patricia Reițman \\
Department of Software Engineering and Automation\\
FCIM, UTM
\end{flushright}

\vspace{3cm}

Chisinau -- 2026

\end{titlepage}

\pagenumbering{arabic}
\setcounter{page}{2}

% ------------------ TABLE OF CONTENTS ------------------
\tableofcontents
\newpage

% ------------------ CHAPTER 1 ------------------
\chapter{INTRODUCTION AND GOAL}

\section{Introduction}
In modern operating systems adhering to the POSIX specification, user applications execute within unprivileged CPU rings (Ring 3 on x86\_64) and are isolated from direct hardware interaction by hardware memory management units (MMUs). Any interaction with peripheral storage, terminal consoles, or IPC pipes requires crossing the hardware privilege boundary into kernel space (Ring 0) via standardized system calls.

While high-level programming interfaces provide abstract wrappers (such as \texttt{printf}, \texttt{fopen}, and \texttt{fread} in the C standard library), these routines ultimately decompose into fundamental kernel system calls (\texttt{openat}, \texttt{write}, \texttt{read}, \texttt{close}). The \texttt{strace} diagnostic utility leverages the Linux kernel's \texttt{ptrace} debugging interface to intercept, decode, and log every system call issued by a process, providing complete transparency into the runtime user-kernel dialogue.

\section{Lab Objective}
The primary objective of this laboratory work is to inspect and analyze the low-level system call interface of Linux using \texttt{strace}. Specifically, the work covers:
\begin{itemize}
    \item tracing the execution of standard GNU utilities (\texttt{ls -l}) and identifying exact system calls responsible for directory examination and terminal output;
    \item developing a standalone C program (\texttt{syscalls.c}) that interacts directly with the kernel using core POSIX file descriptor system calls (\texttt{open}, \texttt{write}, \texttt{read}, \texttt{close});
    \item evaluating the POSIX file descriptor allocation policy and explaining why initial open operations return descriptor \texttt{3};
    \item analyzing the relationship between user-space buffering libraries (\texttt{printf}) and kernel-level output routines (\texttt{write});
    \item profiling and counting system call distributions using the \texttt{strace -c} summary flag, analyzing why file reading requires multiple \texttt{read()} invocations, and observing chunk-based buffer reading and End-Of-File (EOF) detection.
\end{itemize}

% ------------------ CHAPTER 2 ------------------
\chapter{PART 1: MONITORING COMMAND EXECUTION WITH STRACE}

\section{Tracing the \texttt{ls -l} Utility}
The standard directory listing utility \texttt{ls -l} was executed under \texttt{strace}, filtering for system calls associated with descriptor manipulation and data transmission:

\begin{lstlisting}[language=bash]
$ strace -f -e trace=openat,read,write,close ls -l
close(3)                                = 0
openat(AT_FDCWD, "/etc/ld.so.cache", O_RDONLY|O_CLOEXEC) = 3
close(3)                                = 0
openat(AT_FDCWD, "/lib/x86_64-linux-gnu/libc.so.6", O_RDONLY|O_CLOEXEC) = 3
read(3, "\177ELF\2\1\1\3\0\0\0\0\0\0\0\0\3\0>\0\1\0\0\0P\237\2\0\0\0\0\0"..., 832) = 832
close(3)                                = 0
openat(AT_FDCWD, ".", O_RDONLY|O_NONBLOCK|O_CLOEXEC|O_DIRECTORY) = 3
close(3)                                = 0
write(1, "total 28\n-rw-rw-r-- 1 vitalie vi"..., 385) = 385
close(1)                                = 0
close(2)                                = 0
+++ exited with 0 +++
\end{lstlisting}

\begin{figure}[H]
\centering
\caption{Figure 2.1 -- Selective system call trace of ls -l under strace}
\end{figure}

\section{Answers to Required Observation Questions}
\begin{enumerate}
    \item \textbf{System Call Used to Display Text:}
    The \texttt{ls} utility employs the \textbf{\texttt{write}} system call to transmit formatted directory listings to the terminal console:
    \begin{lstlisting}[language=bash]
    write(1, "total 28\ndrwxrwxr-x ...", 385) = 385
    \end{lstlisting}
    Before calling \texttt{write}, \texttt{ls} queries directory metadata using \texttt{openat()} to open the directory file descriptor, \texttt{getdents64()} to enumerate directory entries, and \texttt{newfstatat()} to retrieve file permissions, timestamps, and sizes.

    \item \textbf{File Descriptor Written To:}
    The descriptor number passed to the \texttt{write()} system call is strictly \textbf{\texttt{1}}, representing standard output (\texttt{stdout}).
\end{enumerate}

% ------------------ CHAPTER 3 ------------------
\chapter{PART 2: DIRECT KERNEL CALLS (\texttt{syscalls.c})}

\section{Implementation of \texttt{syscalls.c}}
A C program was constructed to interact directly with the Linux kernel using POSIX system calls (\texttt{open}, \texttt{write}, \texttt{read}, and \texttt{close}) without library-level buffering for file operations:

\begin{lstlisting}[language=C]
/* syscalls.c - open, write, read and close, using only system calls.
 * Build: gcc -Wall -Wextra -o syscalls syscalls.c
 * Run:   ./syscalls
 * See the calls: strace ./syscalls
 */
#include <fcntl.h>  /* open */
#include <unistd.h> /* write, read, close */
#include <string.h> /* strlen */
#include <stdio.h>  /* printf */

int main(void) {
    const char *msg = "hello, kernel\n";

    // open() is a system call: ask kernel for file, get back descriptor
    int fd = open("note.txt", O_CREAT | O_WRONLY | O_TRUNC, 0644);
    if (fd < 0) { 
        perror("open"); 
        return 1; 
    }
    printf("open returned descriptor %d\n", fd);

    // write() is a system call: transfers bytes into kernel buffer cache
    ssize_t n = write(fd, msg, strlen(msg));
    printf("write sent %zd bytes\n", n);
    close(fd);

    // read it back through the kernel
    fd = open("note.txt", O_RDONLY);
    if (fd < 0) { 
        perror("open"); 
        return 1; 
    }
    char buf[64];
    n = read(fd, buf, sizeof(buf) - 1);
    if (n >= 0) {
        buf[n] = '\0';
    }
    printf("read got %zd bytes: %s", n, buf);
    close(fd);

    return 0;
}
\end{lstlisting}

\begin{figure}[H]
\centering
\caption{Figure 3.1 -- Source code of syscalls.c}
\end{figure}

\section{Execution and System Call Trace Verification}
The application was compiled using GCC with strict error checking and executed under \texttt{strace}:

\begin{lstlisting}[language=bash]
$ gcc -Wall -Wextra -o syscalls syscalls.c
$ ./syscalls
open returned descriptor 3
write sent 14 bytes
read got 14 bytes: hello, kernel

$ strace -e trace=openat,write,read,close ./syscalls
openat(AT_FDCWD, "note.txt", O_WRONLY|O_CREAT|O_TRUNC, 0644) = 3
write(1, "open returned descriptor 3\n", 27) = 27
write(3, "hello, kernel\n", 14)           = 14
write(1, "write sent 14 bytes\n", 20)      = 20
close(3)                                = 0
openat(AT_FDCWD, "note.txt", O_RDONLY)  = 3
read(3, "hello, kernel\n", 63)          = 14
write(1, "read got 14 bytes: hello, kernel\n", 33) = 33
close(3)                                = 0
+++ exited with 0 +++
\end{lstlisting}

\begin{figure}[H]
\centering
\caption{Figure 3.2 -- Intercepted system call trace for ./syscalls}
\end{figure}

\section{Answers to Required Observation Questions}
\begin{enumerate}
    \item \textbf{File Descriptor Allocation Mechanics (Why Descriptor 3?):}
    \begin{itemize}
        \item Under POSIX operating systems, every process is initialized with three pre-allocated standard streams:
        \begin{itemize}
            \item \texttt{0} --- Standard Input (\texttt{stdin});
            \item \texttt{1} --- Standard Output (\texttt{stdout});
            \item \texttt{2} --- Standard Error (\texttt{stderr}).
        \end{itemize}
        \item The Linux kernel follows a strict lowest-available-index allocation policy for process file descriptor tables (\texttt{struct files\_struct}). When \texttt{open()} is invoked, the kernel scans the table starting from index 0 and assigns the first free numerical slot.
        \item Since slots 0, 1, and 2 are already assigned to terminal descriptors upon process initialization, the first available index is \textbf{\texttt{3}}.
        \item When \texttt{close(3)} is called, slot 3 is released. Consequently, the subsequent \texttt{open("note.txt", O\_RDONLY)} reallocates slot \textbf{\texttt{3}}.
    \end{itemize}

    \item \textbf{Relationship Between \texttt{printf()} and \texttt{write()}:}
    \begin{itemize}
        \item \texttt{printf()} is a user-space library function implemented in the C runtime library (\texttt{glibc}). It provides string formatting and user-space memory buffering.
        \item \texttt{write()} is the underlying kernel system call that transfers raw bytes across the user-kernel privilege boundary.
        \item As seen in Figure 3.2, calling \texttt{printf("open returned descriptor 3\textbackslash n", fd)} does not execute in kernel space directly; instead, \texttt{glibc} formats the string into a buffer and immediately invokes \texttt{write(1, "open returned descriptor 3\textbackslash n", 27)} to flush the line to descriptor 1.
    \end{itemize}
\end{enumerate}

% ------------------ CHAPTER 4 ------------------
\chapter{PART 3: COUNTING SYSTEM CALLS WITH STRACE (-c)}

\section{Aggregated System Call Profiling}
In addition to step-by-step stream tracing, \texttt{strace} provides a dedicated profiling mode enabled via the \texttt{-c} (or \texttt{--summary-only}) flag. In this mode, \texttt{strace} attaches to the process, intercepts all system calls, measures the execution duration spent within the kernel for each call, and generates an aggregated statistical summary table upon process termination.

To evaluate real file I/O operations, the GNU \texttt{cat} utility was profiled while reading the source file \texttt{syscalls.c}, redirecting output to \texttt{/dev/null} to isolate file reading from console display latency:

\begin{lstlisting}[language=bash]
$ strace -c cat syscalls.c > /dev/null
\end{lstlisting}

\section{Execution and Summary Table Output}
The command executed and generated the following statistical profile:

\begin{lstlisting}
% time     seconds  usecs/call     calls    errors syscall
------ ----------- ----------- --------- --------- ----------------
 25.17    0.000037          18         2           read
 19.05    0.000028           7         4           mmap
 14.29    0.000021          10         2           openat
 12.24    0.000018           6         3           close
  9.52    0.000014           7         2           newfstatat
  6.80    0.000010          10         1           write
  4.76    0.000007           7         1           munmap
  4.08    0.000006           6         1           brk
  4.08    0.000006           6         1           execve
  0.00    0.000000           0         1           arch_prctl
  0.00    0.000000           0         1           set_tid_address
  0.00    0.000000           0         1           set_robust_list
  0.00    0.000000           0         1           prlimit64
  0.00    0.000000           0         1           rseq
------ ----------- ----------- --------- --------- ----------------
100.00    0.000147                    22           total
\end{lstlisting}

\begin{figure}[H]
\centering
\caption{Figure 4.1 -- System call count and timing summary generated by strace -c}
\end{figure}

\section{Answers to Required Observation Questions}
\begin{enumerate}
    \item \textbf{From the summary, how many \texttt{read} calls did \texttt{cat} make?}
    
    From the summary table, \texttt{cat} made exactly \textbf{2} \texttt{read} system calls.

    \item \textbf{Explain in one sentence why reading a file needs more than one \texttt{read}:}
    
    \textbf{A file is read in buffer-sized chunks, and the program cannot know when the end of the file is reached until a subsequent \texttt{read()} system call returns 0 bytes (the standard POSIX EOF indicator), which requires an extra call after all data has been ingested.}
\end{enumerate}

\section{Architectural Analysis: Chunk-Based Ingestion and EOF Detection}
When \texttt{cat} reads a file, it allocates a memory buffer according to the file system's block size (typically queried via \texttt{fstat}, e.g., 64 KiB or 131,072 bytes). 
\begin{itemize}
    \item \textbf{First \texttt{read()} call:} The kernel transfers all available bytes of \texttt{syscalls.c} (984 bytes) from the kernel page cache into the user buffer. The system call returns $984$, representing the number of bytes read.
    \item \textbf{Second \texttt{read()} call:} Although the program received fewer bytes than its buffer capacity, the POSIX specification permits short reads even when more data might follow (for instance, on pipes, network sockets, or slow storage). Therefore, the utility must issue another \texttt{read()} to definitively check for additional bytes. The kernel identifies that the file pointer has reached the end of the file and returns \textbf{0}.
    \item \textbf{EOF Termination:} In Unix/POSIX convention, a return value of \texttt{0} from \texttt{read()} represents the End-Of-File (EOF) condition. Upon receiving 0, \texttt{cat} terminates its read loop and closes the file descriptor.
\end{itemize}

% ------------------ CHAPTER 5 ------------------
\chapter{ANALYSIS OF KERNEL-USER BOUNDARIES}

\section{System Call Dispatching Architecture}
On the x86\_64 Linux architecture, the transition from unprivileged userland to privileged kernel mode proceeds through the \texttt{syscall} machine instruction:
\begin{enumerate}
    \item \textbf{Register Setup:} User-space wrappers load the system call number into register \texttt{rax} (e.g., \texttt{rax = 1} for \texttt{SYS\_write}, \texttt{rax = 257} for \texttt{SYS\_openat}). Call parameters are loaded sequentially into registers \texttt{rdi}, \texttt{rsi}, \texttt{rdx}, \texttt{r10}, \texttt{r8}, and \texttt{r9}.
    \item \textbf{Privilege Transition:} The CPU executes the \texttt{syscall} instruction. The hardware transitions from Ring 3 to Ring 0, saves the instruction pointer into \texttt{rcx}, and jumps to the kernel entry point recorded in the Model Specific Register \texttt{IA32\_LSTAR}.
    \item \textbf{Dispatcher Execution:} The kernel entry stub pushes user registers to the kernel stack, checks the system call number against the architecture jump table (\texttt{sys\_call\_table}), and executes the corresponding service routine (such as \texttt{ksys\_write()}).
    \item \textbf{Return to Userland:} Upon completing the request, the return value or negative error code is placed into \texttt{rax}, and the \texttt{sysretq} instruction restores user-space execution.
\end{enumerate}

\section{The Role of \texttt{strace} and \texttt{ptrace}}
The \texttt{strace} utility attaches to the target process using the \texttt{ptrace(PTRACE\_SYSCALL, ...)} primitive. The kernel suspends the traced process twice per system call: once immediately before syscall entry (allowing inspection of arguments), and once immediately following syscall exit (allowing inspection of return values). This architecture ensures that every transition is captured without altering application source code.

% ------------------ CHAPTER 6 ------------------
\chapter{CONCLUSIONS}
In this laboratory work, the low-level communication interface between user-space applications and the Linux kernel was analyzed using the \texttt{strace} diagnostic utility across three comprehensive parts:
\begin{enumerate}
    \item \textbf{Real Command Monitoring (Part 1):} Executing \texttt{ls -l} under syscall interception demonstrated that standard utilities depend on basic kernel primitives (\texttt{openat}, \texttt{getdents64}, \texttt{write}) for resource discovery and data rendering.
    \item \textbf{Direct Kernel Interaction (Part 2):} Developing \texttt{syscalls.c} confirmed that file descriptor allocations follow a deterministic lowest-index policy starting at descriptor \texttt{3}, and demonstrated that standard formatting functions (\texttt{printf}) are user-space abstractions that delegate directly to \texttt{write()}.
    \item \textbf{System Call Profiling (Part 3):} Profiling \texttt{cat syscalls.c > /dev/null} with \texttt{strace -c} demonstrated quantitative call counting and proved that file ingestion operates in buffer-sized chunks, requiring a terminal 0-byte return to detect the End-Of-File condition.
\end{enumerate}

The experiment highlighted the structural separation of privilege rings, reinforcing the role of the operating system kernel as the sole arbitrator of hardware resources.

% ------------------ REFERENCES ------------------
\chapter*{REFERENCES}
\addcontentsline{toc}{chapter}{REFERENCES}

\noindent [1] Michael Kerrisk, \textit{The Linux Programming Interface: A Linux and UNIX System Programming Handbook}, No Starch Press, 2010.

\noindent [2] Robert Love, \textit{Linux System Programming: Talking Directly to the Kernel and C Library}, O'Reilly Media, 2nd Edition, 2013.

\noindent [3] W. Richard Stevens and Stephen A. Rago, \textit{Advanced Programming in the UNIX Environment}, Addison-Wesley Professional, 3rd Edition, 2013.

\noindent [4] Department of Software Engineering and Automation, \textit{Operating Systems Lab 9: Watching System Calls with strace}, FCIM, Technical University of Moldova, 2026.

\end{document}
